\section{Nomination persistante des brins}
\label{sec:reevaluation-nomination-persistante-brins}

Dans le chapitre~\ref{cha:etat-art-reevaluation}
Section~\ref{sec:etat-art-reevaluation-synthese-cardot}
(page~\pageref{sec:etat-art-reevaluation-synthese-cardot}), nous avons
rapidement résumé les travaux de Cardot sur lesquels nous nous appuyons. %
En particulier, nous avons décrit son système de nomination persistante des
entités topologiques. %

Dans nos travaux, nous avons constaté que chaque opération appliquée sur une orbite implique la création d'un nom persistant, alors l'utilisation de l'historique d'un unique brin est suffisant pour nommer une orbite. %
Toute autre séquence d'évolutions de cette orbite serait, de toutes manières,
représentée par un autre nom persistant et réévalué en fonction de celui-ci. %

Ainsi, dans cette section, nous détaillons notre système de nomination
persistante pour les brins en évaluant pas à pas la construction nº1%


\begin{figure}[htb!]
	\centering
	\begin{subfigure}{.6\textwidth}
		\centering
		\includegraphics[height = 160pt]{rule-create-square-face}
		\caption{}
	\end{subfigure}
	\begin{subfigure}{.35\textwidth}
		\centering
		\includegraphics[width = 100pt]{example-image-b}
		\caption{}
	\end{subfigure}
	\caption{\texttt{1-carre(pos)}}
\end{figure}

\maxime{Dans ce paragraphe, je montre la construction de l'historique pour un
  échantillon de brin et plus particulièrement les brins 3 et 6.}

\begin{figure}[htb!]
	\centering
	\begin{subfigure}{.6\textwidth}
		\centering
		\includegraphics[height = 90pt]{rule-face-extrusion-compact}
		\caption{}
	\end{subfigure}
	\begin{subfigure}{.35\textwidth}
		\centering
		\includegraphics[width = 100pt]{example-image-b}
		\caption{}
	\end{subfigure}
	\caption{\texttt{2-extrusionFace(\pndart{6}{1n6})}}
\end{figure}

\maxime{Dans ce paragraphe, je montre que le paramètre topologique d'une règle
  est une orbite incidente à un brin, donc un nom persistant est une
  photographie de l'historique d'un brin au moment de l'application d'une règle
  à une étape donnée. %
  Je montre la mise à jour des historiques et la construction des PNs suivants.}

\begin{figure}[htb!]
	\centering
	\begin{subfigure}{.6\textwidth}
		\centering
		\includegraphics[height = 40pt]{rule-insert-vertex}
		\caption{}
	\end{subfigure}
	\begin{subfigure}{.35\textwidth}
		\centering
		\includegraphics[width = 100pt]{example-image-b}
		\caption{}
	\end{subfigure}
	\caption{\texttt{3-insertionSommet(\pndart{46}{1n5;2n5})}}
\end{figure}

\begin{figure}[htb!]
	\centering
	\begin{subfigure}{.6\textwidth}
		\centering
		\includegraphics[height = 90pt]{rule-volume-face-extrusion-compact-color}
		\caption{}
	\end{subfigure}
	\begin{subfigure}{.35\textwidth}
		\centering
		\includegraphics[width = 100pt]{example-image-b}
		\caption{}
	\end{subfigure}
	\caption{\texttt{4-extrusionFaceVolume(\pndart{47}{1n5;2n6;3n0})}}
\end{figure}

\begin{figure}[htb!]
	\centering
	\begin{subfigure}{.6\textwidth}
		\centering
		\includegraphics[height = 40pt]{rule-face-triangulation-3D}
	\end{subfigure}
	\begin{subfigure}{.35\textwidth}
		\centering
		\includegraphics[width = 100pt]{example-image-b}
	\end{subfigure}
	\caption{\texttt{6-contractionFace(\pndart{63}{1n4;2n6;4n6})}}
\end{figure}

\begin{figure}[htb!]
	\centering
	\begin{subfigure}{.6\textwidth}
		\centering
		\includegraphics[height = 40pt]{rule-face-collapse}
	\end{subfigure}
	\begin{subfigure}{.35\textwidth}
		\centering
		\includegraphics[width = 100pt]{example-image-b}
	\end{subfigure}
	\caption{Contraction de la face du dessus}
\end{figure}

\begin{figure}[htb!]
	\centering
	\begin{subfigure}{.6\textwidth}
		\centering
		\includegraphics[height = 40pt]{rule-vertex-chamfer-color}
	\end{subfigure}
	\begin{subfigure}{.35\textwidth}
		\centering
		\includegraphics[width = 100pt]{example-image-b}
	\end{subfigure}
	\caption{\texttt{7-chanfrein(\pndart{55}{1n5;2n6;3n0;4n4;6n2})}}
\end{figure}

\begin{figure}[htb!]
	\centering
	\begin{subfigure}{.6\textwidth}
		\centering
		\includegraphics[height = 40pt]{rule-colour-face}
	\end{subfigure}
	\begin{subfigure}{.35\textwidth}
		\centering
		\includegraphics[width = 100pt]{example-image-b}
	\end{subfigure}
	\caption{\texttt{7-coloration(\pndart{35}{1n3;2n4;5n0})}}
\end{figure}

%%% Local Variables:
%%% mode: LaTeX
%%% TeX-master: "../../../main"
%%% End:
